Cieľom tohto zadania bolo vytvoriť simuláciu modelu synchrónneho generátora s použitím modelu budiča ST4C v prostredí Simulink a overiť dynamické vlastnosti modelu budiča pre synchrónny generátor. Dynamické vlastnosti boli overené pomocou simulácie a následného vyhodnotenia prechodových a frekvenčných charakteristík. 

Do schémy modelu bola implementovaná vodná turbína s regulátorom výkonu pomocou bloku HTG. Pre vizualizáciu a porovnanie výsledkov boli vykonané simulácie, pomocou ktorých boli získané odozvy svorkového napätia a činného výkonu na skokové zmeny referenčného napätia. Optimalizáciou parametrov stabilizátora PSS a regulátora turbíny bolo možné dosiahnuť rýchlejšie tlmenie kmitania svorkového napätia a činného výkonu synchrónneho generátora po skokovej zmene referenčného napätia. Optimalizované parametre boli následne použité pre simuláciu, z ktorej boli získané prechodové a frekvenčné charakteristiky, ktoré sú v tejto dokumentácii prezentované a porovnané. Optimalizácia parametrov bola realizovaná pomocou genetického algoritmu implementovaného ako skript v MATLABe, ktorý iteratívne upravoval parametre a vyhodnocoval výsledné prechodové charakteristiky podľa zvolenej kriteriálnej funkcie. Kriteriálnou funkciou optimalizácie bola suma absolútnych hodnôt odchýlok svorkového napätia a činného výkonu od ich referenčných hodnôt.